\section{Conclusions}\label{sec:conclusions}

We proposed the flicker-bounded phase-continuous sequence for moving optical traps, in which every frame is a weighted Gerchberg--Saxton step whose trap phases are locked to a smooth prescribed trajectory, started from the previous hologram and weights, with a frame-count equation that bounds the shift of the pattern between frames, and we measured it against the two published sequences on one simulated model. Against linear phase interpolation the sequence is clearly better: on the acceptance task the lowest moving trap keeps \accBTransMean{} of its mean against \accLpiTransMean{}, with a paired difference of \diffTwoMean{} (95 percent interval \diffTwoCiLo{} to \diffTwoCiHi{}) over \nSeedsTwo{} seeds, and a per-frame spread that is lower by a factor of \fpeval{round(\accLpiSigmaMed/\accBSigmaMed)}. Against per-frame Gerchberg--Saxton with phase induction, the published method that also optimizes every frame, the result is a trade: the proposed sequence has a largest phase step \stepRatio{} times smaller, \accBStepMax{} rad against \accKimFiveStepMax{} rad, and a per-frame spread \sigmaRatio{} times smaller, but a lower transient minimum, \accBTransMean{} against \accKimFiveTransMean{}, a tie at equal work and at 25 frames, and a cost of \secRatio{} times the seconds at its own iteration count.

On the long moves, with the frames that the bound requires, the sequence reaches $0.9$ on \longAMet{} of \longSeeds{} seeds of each lattice move and on \cTenMet{} of \cTenSeeds{} seeds of the sort, where the published per-frame method reaches it on every seed of the lattice moves; the failing seeds are not rescued by doubling the frames, which points to the prescribed phase path. The ablation shows that the per-frame optimization is necessary, that the lock lowers the phase step by a factor of \lockStepRatio{} and narrows the spread of outcomes without being what keeps the mean above $0.9$, that three iterations per frame are enough on the acceptance task, and that a linear time profile gives a higher dip than the minimal-jerk profile.

The novelty that survives the comparison is the prescribed and locked phase trajectory with a derived frame-count bound, and what it buys is a controlled phase step, not a higher dip. The next steps are to measure the sequences on the atoms, where the survival against the phase step settles whether the smaller step has a use; to repeat the comparison with a shared time profile, so that the dip is attributed to the frame computation alone; to compare at equal wall-clock time on the long moves; and to replace the calibrated step bound by one derived for several traps.
